\documentclass[10pt,twocolumn,letterpaper]{article}
\usepackage[rebuttal,applications]{wacv}  % use this for an Applications Track rebuttal
%\usepackage[rebuttal,algorithms]{wacv}  % use this for an Algorithms Track rebuttal

% Include other packages here, before hyperref.
\usepackage{graphicx}
\usepackage{amsmath}
\usepackage{amssymb}
\usepackage{booktabs}

% Import additional packages in the preamble file, before hyperref
\input{preamble}

% If you comment hyperref and then uncomment it, you should delete
% egpaper.aux before re-running latex.  (Or just hit 'q' on the first latex
% run, let it finish, and you should be clear).
\definecolor{wacvblue}{rgb}{0.21,0.49,0.74}
\usepackage[pagebackref,breaklinks,colorlinks,allcolors=wacvblue]{hyperref}
\usepackage{xcolor}
\newcommand{\rBct}{\textcolor{wacvblue}{rBct}}
\newcommand{\uizG}{\textcolor{orange}{uizG}}
\newcommand{\wrik}{\textcolor{teal}{wrik}}
\newcommand{\AC}{\textcolor{red}{AC}}

% If you wish to avoid re-using figure, table, and equation numbers from
% the main paper, please uncomment the following and change the numbers
% appropriately.
%\setcounter{figure}{2}
%\setcounter{table}{1}
%\setcounter{equation}{2}

% If you wish to avoid re-using reference numbers from the main paper,
% please uncomment the following and change the counter value to the
% number of references you have in the main paper (here, 100).
%\makeatletter
%\apptocmd{\thebibliography}{\global\c@NAT@ctr 100\relax}{}{}
%\makeatother

%%%%%%%%% PAPER ID  - PLEASE UPDATE
\def\wacvPaperID{967} % *** Enter the WACV Paper ID here
\def\confName{WACV}
\def\confYear{2026}

\begin{document}

%%%%%%%%% TITLE - PLEASE UPDATE
\title{Manga109Script Dataset for Script to Layout Generation }  % **** Enter the paper title here

\maketitle
\thispagestyle{empty}
\appendix


We thank the area chair and reviewers for their thoughtful feedback. 
We are encouraged that they recognized the value of releasing the \textit{Manga109Script} dataset with human-validated annotations and benchmark (\AC, \rBct, \uizG, \wrik), the novelty of defining the Script2Layout task and the clarity of our methodology with extensive experiments (\uizG, \wrik), and the practical significance of bridging narrative scripts and visual layouts (\AC, \wrik). We address reviewer comments below and will incorporate all feedback.

\noindent\textbf{@\wrik, \rBct – Concerns on the limitations of Manga109 (age, diversity, and color applicability):}
Manga109 has been recognized as a standard benchmark dataset for manga analysis in recent top-tier venues (e.g., MangaUB, IEEE Multimedia ’25; Advancing Manga Analysis, CVPR ’25), and its diversity in style and content has been acknowledged. Moreover, prior studies show that color images often facilitate recognition and understanding (The Effect of Color Space Selection on Detectability and Discriminability of Colored Objects ’17), suggesting that our script construction method (Sec.\,3.1) is also applicable to other datasets, including those with color comics. Therefore, we expect that the effectiveness of Script2Layout training should not fundamentally depend on whether the data extracted from color or grayscale comics.


\noindent\textbf{@\uizG – GPT-4.1 is enough to recognize the dialogue order.}
Regarding this comment, we conducted an additional experiment by providing GPT-4.1 with 109 pages (one page per work) from Manga109 (Sec.,3.1). The correct dialogue order was identified in only 26 out of 109 pages. In contrast, applying the algorithm (Magi ’24) used in this work to the same subset achieved 96 out of 109, demonstrating the effectiveness of our script construction method. 

\noindent\textbf{@\uizG – Evaluation metrics are questionable:}  
We agree that layout generation inherently admits multiple valid solutions, so comparison with a single ground truth is insufficient. However, prior work has also noted the difficulty of ensuring reliability in human evaluation (LTSim ’24). Furthermore, given the limited accuracy of GPT-4.1 on dialogue order (Sec.\,3.1), we question whether MLLM-based evaluation is fully reliable in this setting. To address these concerns, we introduced two additional metrics: Overlap of Speech Balloons (Ovp.) and LTSim-MMD (Sec.\,4.3). Ovp. captures universally intuitive penalties for readability, while LTSim-MMD measures geometric similarity in a many-to-many manner, making it suitable for evaluating the plausibility of manga layouts. Under these metrics as well, the value of using scripts for layout generation is confirmed (Sec.\,4.5).

\noindent\textbf{@\uizG – Lack of innovative methods and only small improvements from fine-tuning:}  
We would like to note that, in the Applications Track, methodological innovativeness is not a strict requirement. Moreover, our study demonstrates consistent effects of scripts across multiple metrics and LLMs, indicating that the improvements are meaningful. There is no established criterion for what constitutes a “small” improvement.

\noindent\textbf{@\uizG – Training and test sets are too similar:}  
The results of our experiments, conducted with a split at the work level (Sec.,4.2), should not be regarded as suffering from overly high similarity between training and test data, because of Manga109's diversity.

\noindent\textbf{@\wrik – Lack of technical innovation, especially handling multiple pages in script construction:}  
Our script generation method using summaries (Sec.\,3.1) explicitly addresses understanding across multiple pages. In addition, we would like to clarify that, in the Applications Track, methodological innovativeness is not a mandatory requirement.

\noindent\textbf{@\uizG – Existing methods also require layout elements as input, so the difference from Script2Layout is unclear:}
While layout elements are not technically a mandatory requirement for the Script2Layout task, we consider them practically necessary to ensure controllability in real-world applications. We do not aim to develop a full-automatic system, but an interactive system.


\noindent\textbf{@\wrik – Script2Layout is unclear how this task contributes to the final generation of manga.}
The design of the Script2Layout task is grounded in the actual practice of creating manga from narratives. To avoid the rework cost that can arise when manga images are generated directly from stories, we propose generating layouts as an intermediate representation. This allows controllability in image generation, similar to how ControlNet enables control, and we argue that this intermediate step is highly beneficial for practical manga creation.

\noindent\textbf{@\uizG, \wrik – Concerns on the narrow scope of utility of Manga109DScript.}
It can be said that the true significance of our dataset lies in the fact that it has inspired a sense of various possibilities.

\noindent\textbf{@\rBct – Lack of detailed analysis of T5 vs. DeepSeek performance differences:}  
We agree that analyzing performance differences across architectures is an important research direction. However, such an in-depth study is beyond the scope of this Applications Track and not a mandatory requirement for acceptance. Instead, our focus is on demonstrating that the constructed scripts bring consistent benefits across different models, thereby underscoring the significance of Script2Layout as a task.

\noindent\textbf{@\uizG, \wrik, \rBct – Errors in the description and figures that are hard to read}
Thank you for pointing that out. We have addressed all of these issues. (Sec.\,3.1, 4.1, Fig.\,1-3)



\end{document}
